\section{Product-of-Experts Decoding Framework}

\label{sec:poe}

\subsection{Token-Level PoE Formulation}

At each decoding step $t$, the adapted distribution is a PoE of the base model
and $M$ experiential priors:
\begin{equation}
    \pi_{\text{ASO}}(y_t \mid x, y_{<t})
    = \frac{1}{Z_t} \; \pi_\theta(y_t \mid x, y_{<t})
    \prod_{m=1}^{M} \phi_m(y_t \mid x, y_{<t})^{\eta_m},
    \label{eq:poe-main}
\end{equation}
where:
\begin{itemize}[leftmargin=1.5em]
    \item $\pi_\theta$ is the frozen base model (expert 0),
    \item $\phi_m$ is the $m$-th experiential prior factor,
    \item $\eta_m \geq 0$ is the weight for prior $m$,
    \item $Z_t = \sum_{v \in \mathcal{V}} \pi_\theta(v \mid x, y_{<t})
          \prod_m \phi_m(v \mid x, y_{<t})^{\eta_m}$ is the normalizing constant.
\end{itemize}

In log-space, the PoE decoding reduces to a simple addition of logits:
\begin{equation}
    \log \pi_{\text{ASO}}(y_t = v \mid x, y_{<t})
    = \ell_{t,v} + \sum_{m=1}^{M} \eta_m \log \phi_m(v \mid x, y_{<t}) - \log Z_t,
    \label{eq:poe-logspace}
\end{equation}
where $\ell_{t,v}$ are the base model logits. This is computationally
inexpensive: the overhead is $O(MV)$ additions per decoding step, negligible
compared to the $O(d^2)$ cost of the transformer forward pass.

\subsection{Expert Factor Construction}

Each experiential prior $\phi_m$ is constructed from the semantic advantage
$S_m$ extracted during TF-GRPO. The construction proceeds in three steps:

\paragraph{Step 1: Token Set Identification.}
Given pattern $s_k = (\text{pattern}_k, c_k, \mathcal{V}_k)$, identify the set
of vocabulary tokens $\mathcal{V}_k$ that are promoted or demoted by this
pattern. For code generation, this typically involves:
\begin{itemize}[leftmargin=1.5em]
    \item Function/method names (e.g., \texttt{get} vs.\ \texttt{\_\_getitem\_\_}),
    \item Control flow tokens (e.g., \texttt{try}/\texttt{except} vs.\ \texttt{if}/\texttt{else}),
    \item Library-specific tokens (e.g., \texttt{pandas} operations).
\end{itemize}

\paragraph{Step 2: Delta Computation.}
For each token $v \in \mathcal{V}_k$, compute the log-probability adjustment:
\begin{equation}
    \Delta_m(v) = c_m \cdot \begin{cases}
        +\delta_+ & \text{if } v \text{ is promoted by pattern } m, \\
        -\delta_- & \text{if } v \text{ is demoted by pattern } m, \\
        0 & \text{otherwise},
    \end{cases}
\end{equation}
where $\delta_+, \delta_- > 0$ are the promotion/demotion magnitudes (default
$\delta_+ = \delta_- = 1.0$). The confidence $c_m$ scales the adjustment.

\paragraph{Step 3: Factor Assembly.}
The expert factor is:
\begin{equation}
    \log \phi_m(v \mid x, y_{<t}) = \begin{cases}
        \Delta_m(v) & \text{if context matches pattern } m \text{ and } v \in \mathcal{V}_m, \\
        0 & \text{otherwise}.
    \end{cases}
\end{equation}
The context matching condition ensures that priors are applied only when
relevant. For example, a pattern about error handling is applied only when the
model is generating code within a function body that performs I/O operations.

\subsection{Expert Weight Derivation}

The weight $\eta_m$ for prior $m$ is derived from its quality score $q_m \in (0, 1)$
via a Bayesian argument.

\begin{theorem}[Optimal Expert Weights]
\label{thm:weights}
Under the assumption that each prior $m$ is an independent binary classifier of
``good'' vs.\ ``bad'' token choices with accuracy $q_m$, the Bayes-optimal
log-weight is:
\begin{equation}
    \eta_m = \log \frac{q_m}{1 - q_m}.
    \label{eq:eta-bayes}
\end{equation}
\end{theorem}

\begin{proof}
Consider token $v$ at position $t$. Let $H_+$ denote the event that $v$ is a
good choice and $H_-$ that it is bad. Prior $m$ provides a binary
recommendation: $\phi_m(v) = q_m$ if recommending $v$ (event $R_+$) and
$\phi_m(v) = 1 - q_m$ otherwise. By independence:
\begin{align}
    \frac{P(H_+ \mid R_+^1, \ldots, R_+^M)}{P(H_- \mid R_+^1, \ldots, R_+^M)}
    &= \frac{P(H_+)}{P(H_-)} \prod_{m=1}^M \frac{P(R_+^m \mid H_+)}{P(R_+^m \mid H_-)} \\
    &= \frac{P(H_+)}{P(H_-)} \prod_{m=1}^M \frac{q_m}{1 - q_m}.
\end{align}
Taking logarithms and identifying the base model logit ratio with
$\log(P(H_+)/P(H_-))$, each prior contributes an additive
$\log(q_m / (1 - q_m))$ to the log-odds, which is exactly the PoE with
$\eta_m$ as given.
\end{proof}

\begin{corollary}
Priors with $q_m = 0.5$ (random) receive weight $\eta_m = 0$ and are
automatically ignored. Priors with $q_m > 0.5$ receive positive weight
(contribute constructively), while $q_m < 0.5$ would receive negative weight
(inverted signal). In practice, we clamp $\eta_m \geq 0$ and discard priors
with $q_m < 0.5$.
\end{corollary}

\subsection{Sparse Expert Factors}

In practice, most expert factors are extremely sparse: only a small subset
$\mathcal{V}_m \subset \mathcal{V}$ of tokens are affected by any given pattern.
Let $s_m = |\mathcal{V}_m| / V$ denote the sparsity ratio. Typical values:

\begin{table}[h]
\centering
\begin{tabular}{lcc}
\toprule
Pattern Type & $|\mathcal{V}_m|$ & Sparsity $s_m$ \\
\midrule
API preference & 5--20 & $10^{-4}$ \\
Control flow & 10--50 & $10^{-3}$ \\
Style convention & 20--100 & $10^{-3}$ \\
Library selection & 50--200 & $10^{-2}$ \\
\bottomrule
\end{tabular}
\caption{Typical sparsity of expert factors for code generation tasks.}
\label{tab:sparsity}
\end{table}

The sparse representation enables efficient storage and retrieval. Only non-zero
entries need to be stored, reducing memory from $O(V)$ to $O(|\mathcal{V}_m|)$
per factor.

\subsection{Temperature and Sharpness Control}

The PoE can be sharpened or softened via a global temperature $\tau$:
\begin{equation}
    \pi_{\text{ASO}}^\tau(y_t = v \mid x, y_{<t})
    = \frac{1}{Z_t^\tau} \left[\pi_\theta(v \mid x, y_{<t})
    \prod_{m=1}^{M} \phi_m(v \mid x, y_{<t})^{\eta_m}\right]^{1/\tau}.
\end{equation}
Low $\tau$ sharpens the distribution (more deterministic), while high $\tau$
flattens it (more exploratory). We use $\tau = 0.7$ for final generation and
$\tau = 1.2$ during TF-GRPO exploration phases.

% ============================================================================
% 5. 1-BIT PRIOR COMPRESSION
% ============================================================================
